%!TEX program = lualatex
\documentclass[ngerman,headheight=70pt]{scrartcl}
\usepackage[ngerman]{babel}
\usepackage{blindtext}
\usepackage[utf8]{luainputenc}
\usepackage{scrlayer-scrpage}
\usepackage[T1]{fontenc}
\usepackage{fontspec}
\usepackage{multicol}
\usepackage{enumitem}
\usepackage[german=quotes]{csquotes}
\usepackage[bottom=2cm,footskip=8mm]{geometry}
\setmainfont{Source Sans Pro}
\MakeOuterQuote{"}

\setlength{\parindent}{0pt}
\setlength{\parskip}{6pt}

\pagestyle{plain.scrheadings}
\thispagestyle{scrheadings}
\lohead{\vspace*{-1.5cm}\includegraphics[scale=0.4]{../images/stupa_logo_rot.pdf}}
\rohead{\\Präsidium des\\Studierendenparlaments}

\lofoot[Universität Hamburg · Präsidium des Studierendenparlaments\\\hrule Von-Melle-Park 5 · D-20146 Hamburg]{Universität Hamburg · Präsidium des Studierendenparlaments\\\hrule Von-Melle-Park 5 · D-20146 Hamburg}
\cofoot[]{}
\cohead[\thepage]{}

\renewcommand{\thesubsection}{TOP \arabic{subsection}}
\renewcommand{\thesubsubsection}{\arabic{subsubsection}.}
\setcounter{subsection}{-1}

\begin{document}
    UHH · StuPa-Präsidium · Von-Melle-Park 5 · D-20146 Hamburg

    %%% PAGE ONE (EDIT LATER) %%%%
    % Hier den Titel des Protokolls eintragen
    \section*{Protokoll der 11. Sitzung des Studierendenparlaments vom 1. Dezember 2016}

    % Protokollanten eintragen
    \textbf{Protokoll: Gunhild Berdal}\\
    \textbf{Ort: VMP9 S08}\\
    \textbf{Beginn: 18.45 Uhr}\\
    \textbf{Ende: 21.30 Uhr}

    % Wohl die komplizierteste Stelle im Protokoll
    % Als Faustregel sollten nicht mehr als 3 Namen auf einer Zeile stehen
    % Einfach an dieser Vorlage orientieren
    % Im Idealfall ist hier nicht viel zu ändern, da kaum jemand fehlt.
    % Rücktritte müssen natürlich pro Protokoll angepasst werden
    \vspace{0.5cm}
    \begin{tabular}{ll}
        Anwesend: & \\
            RCDS (5 Sitze): & Kolja Kolb, Ramin Shakiba \\
             CampusGrün (14 Sitze): & Laura Franzen, Geoffrey Youett, Philipp Droll, \\
                                   & Yasemin Günther, Melf Johannsen, Mario Moldenhauer, \\
                                   & Christian Lagod, Jim Martens, Svenja Horn, \\
                                   & Martin Sievert, Armin Günther \\
             BGZM-Liste (2 Sitze) : & \\
             WiWi (2 Sitze): & \\
             Unicorns (5 Sitze): &  Johannes Peplow, Andreas Hartkamp \\
             Liste LINKS (3 Sitze): & Gunhild Berdal, Till Petersen, Sinah Mielich \\
             HWP (2 Sitze): &  \\
             MIN (4 Sitze): &  \\
             SDS* (3 Sitze): & Mena Winkler, Jacob Petersein, Artur Brückmann \\
             Bart-LISTE (2 Sitze): & Timo Zeimet \\
             LHG (1 Sitz): &  \\
             harte zeiten (1 Sitz): & Tobias Berking \\
             Jura (1 Sitz): & \\
             AL (2 Sitze): & Karima Schulze, Henri Weber \\
            & \\
        Entschuldigt: & Ramon Weilinger, Benjamin Welling (RCDS), Elena Rysikova, \\
                      & Ines Friedrich, Dean Haltmeyer (CG), Jakob Pape (BGZM), \\
                      & Charleen Lorenz (WiWi), Annkathrin Löffler, Marielle Hermstrüwer (Unicorns), \\
                      & Ailina Salten (MIN), Daniel Runge (BART)\\
                      &\\
        Unentschuldigt abwesend: & Jennifer Maack (RCDS), Karen Martirosian (BGZM), Claas-Friso Hente (WiWi), \\
                                & Katharina Kucza, Hakan Tasdemir, Gökce Celik (HWP),\\
                                & Lotte Rullkötter, Nabil Jabareen, Jan Detampel (MIN), \\
                                & Tobias Heisig (LHG), Maurice Martens (Jura) \\
                                &\\
        Rücktritte: & Antonia Niecke \rightarrow Kolja Kolb (RCDS) \\
                    & Tahnee Herzig \rightarrow Malte \rightarrow Ines Friedrich (CG) \\
                    & Dominic Laumer \rightarrow Daniel Runge (BART) \\
        Mandatsverluste: &  \\
    \end{tabular}
    \newpage
    %%% Vorgeschlagene Tagesordnung %%%
    \underline{Vorgeschlagene Tagesordnung}
    \begin{enumerate}[label={\textbf{Top \theenumi}},leftmargin=*]
        \item Aktueller Stand StuPa-Wahl (V1617-045) (30 Min.)
        \item Tag gegen Rechts (V1617-043, V1617-044) (45 Min.)
        \item Verschiedenes (5 Min.)
    \end{enumerate}

    \newpage

    % Hier beginnt Teil A. Bei Fortsetzungen kommt hier nur der Teil über
    % die Feststellung der endgültigen Fassung der Tagesordnung.
    % Für sonstige Sitzungen ist dies bereits das Gerüst für die Formalia.
    % Man muss nur noch die protokollierten Inhalte eintragen.

    {\Large \underline{Teil A}}

    \subsection{Formalia}

    \textit{Jim eröffnet die Sitzung um 18.45 Uhr.}

    \subsubsection{Geschäftsbericht Präsidium}

    Jim berichtet:

    Die kleinen Listenvorstellungen sind angekommen und die Briefwahlunterlagen
    rechtzeitig in den Druck gegangen (42 000 Wähler*innen).

    Über die Foyer-Situation (Einschränkung durch Brandschutzbestimmungen) hat
    das Präsidium mit Herrn Krüß von der Uni-Verwaltung gesprochen, der die
    WiWi-Mensa als Ort für den "Markt der Möglichkeiten" vor der
    Listenvorstellung vorgeschlagen hat.

    Mit den Bewerber*innen für die Position als Wahlhelfer*in gab es ein
    produktives Gespräch, ausgewählt sind: Vincent Orth, Ina Reis, Matthias
    Kruspe, Charleen Lorenz und Niklas A. Reininghaus-Zierfuß.

    Die große Wahlzeitung geht nächste Woche in den Druck und soll bis 15.
    Dezember fertig sein.

    \textbf{Rücktritte: Dominic Laumer (Bart-Liste), Daniel Runge rückt nach mit
    Direktmandat. Tahnee Herzig (CampusGrün) tritt zurück, Ines Friedrich (Platz
    21) rückt nach.}

    \subsubsection{Anfragen an das Präsidium}

    Frage Till: Was steht in der Wahlordnung zu Rücktritten und Direktmandaten?
    Bei wie vielen Direktstimmen gilt das?

    Ein Vergleich mit der Wahlordnung bestätigt, dass beim Rücktritt von
    Direktmandatsträger*innen nach absteigender Anzahl von Direktstimmen
    nachgerückt wird.

    Till problematisiert, dass nach dieser Regelung eine Direktstimme mehr zähle
    als die Reihenfolge der Kandidierenden auf der Liste.

    Andreas fragt, wie eigentlich der Stand der Foyer-Nutzung ist.

    Antwort: Im Moment können die Foyers nicht für Veranstaltungen genutzt werden,
    aber die Uni-Verwaltung ist sehr bemüht, in Einzelfällen pragmatische
    Lösungen zu finden.

    Franziska ergänzt: Seit Jahr 2000 gibt es keine Nutzungsbestimmung der
    Foyers. Es soll zuerst ESA O auf Stand gebracht werden, vielleicht bis SoSe
    2017, danach soll es mit den anderen Bereichen zügig gehen.

    Andreas: Gibt es weitere Einschränkungen? In anderen Räumen, bei Kongressen
    o.Ä.?

    Dem Präsidium ist nichts dergleichen bekannt.

    Till: Die politischen Gegner hätten verstanden, dass an der Uni was los ist,
    das könnte dazu beitragen, dass die Mühlen langsam mahlen. In der Feuerwehr
    gebe es ein Erbe aus der Schill-Zeit, das zur Restriktivität gegenüber der
    Uni führen könne. In allen Gremien der Uni-Verwaltung ist man sich einig,
    dass es anders werden muss.

    Es gibt Zustimmung zum Vorschlag, die Listen-Veranstaltung am 15. Dezember
    in der Mensa durchzuführen, und zwar ab 16 Uhr.

    \subsubsection{Geschäftsbericht AStA}

    Es wird vom Studienreformtag berichtet, dass es eine intensive inhaltliche
    Tagung gewesen sei.

    Beim G20-Treffen habe es eine gute gesellschaftliche, historische Diskussion
    gegeben, anknüpfend an frühere Kämpfe. Es soll eine inhaltliche Plattform
    (Manifest gegen G20), eine AStA-Broschüre sowie eine Unterschriftenliste
    geben.

    Aus dem Öffentlichkeitsreferat wird berichtet, dass es eine neue AStA-Homepage
    gibt.

    \subsubsection{Anfragen an den AStA}

    Till: Wie soll die weitere Kommunikation mit dem Präsidium bezüglich der Dies
    Academici laufen?

    Franziska: Im Ausschuss für Lehre und Studium sind zwei Tagungen besprochen
    worden: Eine Arbeitstagung zur Rahmenprüfungsordnung und später ein Dies
    Academicus zur Inklusiven Hochschule.

    Es entspannt sich ein kurzer Austausch über den Zwischenstand der
    Rahmenprüfungsordnung: Die Diskussion sei verkeilt. Es würden reaktionäre
    Positionen über die Studierenden und ihre vermeintlichen sozialen Vorteile
    verbreitet. Es wird positiv bewertet, dass darüber überhaupt diskutiert wird
    und betont, dass die Rahmenprüfungsordnung tatsächlich besser werden müsse.


    \subsubsection{Dringlichkeitsanträge des AStA}

    Keine.

    \subsubsection{Aktuelle Stunde}

    Es liegt kein Antrag vor.

    \subsubsection{Feststellung der endgültigen Fassung des Teils B der Tagesordnung}

    Es gibt keinen Widerspruch gegen die vorgeschlagene Tagesordnung.

    \newpage
    %%% BESCHLOSSENE TAGESORDNUNG %%%
    \underline{Endgültige Fassung der Tagesordnung}
    \begin{enumerate}[label={\textbf{Top \theenumi}},leftmargin=*]
        \item Aktueller Stand StuPa-Wahl (V1617-045) (30 Min.)
        \item Tag gegen Rechts (V1617-043, V1617-044) (45 Min.)
        \item Verschiedenes (5 Min.)
    \end{enumerate}

    \subsubsection{Feststellung der Beschlussfähigkeit}

    \textit{Sieben Minuten Fraktionspause Liste LINKS.}

    \textit{Rücktritte: Antonia Niecke, Emilia Luise Oppermann, Kevin Knackstedt (alle
    RCDS) – Kolja Kolb rückt nach.}

    Mit 25 Parlamentarier*innen ist das Parlament beschlussfähig.


    \subsubsection{Genehmigung der Protokolle der vorangegangenen Sitzungen}

    Das Protokoll der 9. Sitzung wird beschlossen.

    \vspace{0.5cm}
    {\Large \underline{Teil B}}

    \subsection{Aktueller Stand StuPa-Wahl}

    Gunhild stellt den gemeinsamen Antrag des Präsidium zum "Fairen Wahlkampf"
    vor.

    Till fordert alle dazu auf, mit dem Material anderer Listen auch dann pfleglich
    umzugehen, wenn man mit dem Inhalt nicht einverstanden ist. Alle mögen
    politisch mit eigenen Positionen in die Öffentlichkeit gehen und gemäß den
    Grundsätze des Grundgesetzes und des Parteiengesetzes Meinungsbildung
    betreiben.

    Ramin spricht sich für Waffeln und Glühwein aus, um auf Themen aufmerksam zu
    machen und die  Wahlbeteiligung zu erhöhen. Es wird ausführlich pro und
    contra Waffeln diskutiert.

    Als Konsens kann gelten, dass Waffeln nicht alles sind und der Wahlkampf
    politisch sein möge.

    Der Antrag wird bei 2 Enthaltungen einstimmig angenommen.

    \subsection{Tag gegen Rechts}

    Till stellt den Antrag "Nach wie vor: Aufklärung, Kultur, Solidarität!
    Aktionstag gegen Rechts und Veranstaltung am 30. Januar 2017" vor: In
    Auswertung der Trump-Wahl sei ein Tag gegen Rechts sehr sinnvoll. Es gehe
    darum, das wertvolle anti-faschistische Erbe aktuell zur Geltung zu bringen.
    Auch wenn ihre Umfragewerte fallen, die AfD bleibe ein Problem. Der Tag gegen
    Rechts möge stattfinden, Inhalt und Struktur des Antrags soll Unterstützung
    für den Ausschuss sein.

    Der 30. Januar wird mittlerweile knapp, vielleicht Anfang des Sommersemesters
    möglich? Zum 30. Januar möge es eine Veranstaltung des StuPa zum Jahrestag
    der Machtübertragung geben.

    Aus dem Ausschuss gegen wird der Antrag begrüßt. Ein Tag gegen Rechts im SoSe
    sei sinnvoll, auch  wegen der Nähe zur Bundestagswahl.

    Es gibt einige kritische Anmerkungen zum Text, u.a. zum Begriff
    "Superreiche". Richtiger sei, von einer "Konzentration des gesellschaftlichen
    Reichtums in wenigen Händen..." zu sprechen. Es werden entsprechende
    Änderungsanträge gestellt, die von den Antragsstellern übernommen werden.

    Der Gesamtantrag wird mit den Änderungen mit m:1:5 angenommen.

    \subsection{Verschiedenes}

    Der AStA schlägt eine Vollversammlung gegen den Hamburger Haushalt und für
    die Finanzierung der Hochschulen noch im Dezember vor. Nach der Diskussion
    wird vereinbart, am 13. Dezember ein Aktiventreffen zu organisieren, um zu
    klären, wie man in der Stadt zu mehr Bewegung kommt.

    Gunhild berichtet, dass die CDU auf die StuPa-Beschlüsse zum Haushalt
    geantwortet hat und mit allem, was darin gefordert wird, einverstanden sei.

    Am 11.12. soll ein anti-rassistisches Fußballturnier von DIDF stattfinden.

    \textit{Die Sitzung wird um 21.30 Uhr geschlossen.}

\end{document}
